\documentclass[aspectratio=169,10pt]{beamer}
\usetheme{Madrid}
\usecolortheme{default}
\usepackage[fontset=fandol]{ctex}
\usepackage{amsmath,amssymb,booktabs,mathtools,graphicx,hyperref,tikz,array}
\usetikzlibrary{arrows.meta,positioning,decorations.pathreplacing,fit}
\setbeamertemplate{navigation symbols}{}
\setbeamertemplate{footline}[frame number]
\setbeamerfont{frametitle}{size=\large}
\definecolor{positive}{HTML}{0173B2}
\definecolor{negative}{HTML}{DE8F05}
\definecolor{emphasis}{HTML}{029E73}
\definecolor{neutral}{gray}{0.55}
\definecolor{cbPurple}{HTML}{CC78BC}
\newcommand{\pos}[1]{\textcolor{positive}{#1}}
\newcommand{\con}[1]{\textcolor{negative}{#1}}
\newcommand{\HL}[1]{\textcolor{emphasis}{#1}}
\newcommand{\mute}[1]{\textcolor{neutral}{#1}}
\tikzset{
  box/.style={draw, rounded corners, minimum height=0.8cm, font=\small, align=center, fill=positive!10, text width=2.3cm},
  headbox/.style={draw, rounded corners, minimum height=0.7cm, font=\small, align=center, fill=emphasis!12, text width=1.9cm},
  arr/.style={-{Stealth}, thick},
}

\title[awesome\_starvla]{从 StarVLA 到 VLAct：VLA 持续预训练的表示中心路线}
\subtitle{三篇论文 + 一个代码库的系统调研}
\author{awesome\_starvla 调研汇报}
\institute{github.com/asimfish/awesome\_starvla}
\date{2026 年 9 月}

\begin{document}

\begin{frame}
  \titlepage
\end{frame}

% ---------------------------------------------------------------- 2
\begin{frame}{开场：同一个骨干，预训练是双刃剑还是净收益？}
  \textbf{同样是 Qwen3-VL-4B，同样含 InternData-A1 的机器人数据，两篇论文给出相反的结论。}
  \vspace{4pt}
  \begin{center}\footnotesize
  \begin{tabular}{p{2.5cm}p{3.6cm}p{5.4cm}p{1.2cm}}
    \toprule
    论文 & 预训练方式 & RoboCasa-GR1（预训练未见的人形） & 结论 \\
    \midrule
    StarVLA-$\alpha$（2026-04） & 朴素动作拟合 + OXE & 24$\times$10 demos：9.8 $\to$ \con{1.2} & 双刃剑 \\
    StarVLA-$\alpha$（2026-04） & 朴素动作拟合 + InternData-A1 & 24$\times$1000：53.8 $\to$ \con{35.4} & 双刃剑 \\
    VLAct（2026-08） & 表示中心配方 & 全量 48.8 $\to$ \HL{54.0}；20\% 数据 \HL{49.5}，已超全量 GR00T-N1.6（47.6） & 净收益 \\
    \bottomrule
  \end{tabular}
  \end{center}
  \vspace{4pt}
  \begin{itemize}
    \item 差别不在数据量，在\textbf{轨迹如何塑造骨干表征}
    \item 本次汇报：StarVLA 代码库（基础设施）$\to$ StarVLA-$\alpha$（对照板）$\to$ VLAct（配方）$\to$ 下一步
  \end{itemize}
\end{frame}

% ---------------------------------------------------------------- 3
\begin{frame}{三篇论文的关系}
  \begin{center}
  \begin{tikzpicture}[node distance=0.6cm]
    \node[box, text width=3.4cm] (code) {\textbf{StarVLA 代码库}\\ arXiv 2604.05014\\ 骨干–动作头两条契约\\ 4 头 $\times$ 2 类骨干 $\times$ 13 基准};
    \node[box, text width=3.4cm, right=1.2cm of code] (alpha) {\textbf{StarVLA-$\alpha$}\\ arXiv 2604.11757（ECCV 26）\\ 最简基线 = 强 VLM + MLP\\ 重新检验三个"常识"};
    \node[box, text width=3.4cm, right=1.2cm of alpha, fill=emphasis!12] (vlact) {\textbf{VLAct}\\ arXiv 2608.27550\\ 表示中心持续预训练\\ 冻浅层 + 三头 + 统一动作空间};
    \draw[arr] (code) -- node[above, font=\scriptsize] {提供基础设施} (alpha);
    \draw[arr] (alpha) -- node[above, font=\scriptsize] {"预训练是双刃剑"} node[below, font=\scriptsize] {$\Rightarrow$ 为什么？} (vlact);
  \end{tikzpicture}
  \end{center}
  \vspace{6pt}
  \begin{itemize}
    \item 共同点：同一团队（HKUST / CUHK / SmartMore，Jiaya Jia 组）、同一骨干、同一代码
    \item 这让\textbf{动作头、预训练配方、动作空间}三个变量第一次可以被单独隔离
    \item 本仓库：3 篇论文中英 PDF、7 份精读报告、~100 篇文献编目、本 PPT
  \end{itemize}
\end{frame}

% ---------------------------------------------------------------- 4
\section{StarVLA 代码库}
\begin{frame}{StarVLA：两条契约把 VLA 拆成可换部件}
  \begin{center}
  \begin{tikzpicture}[node distance=0.5cm]
    \node[box, text width=1.6cm, fill=neutral!15] (obs) {原始观测\\ \scriptsize RGB 多视角 + 指令 (+ 状态)};
    \node[box, text width=2.3cm, right=0.5cm of obs] (vlm) {VL 骨干\\ \scriptsize Qwen2.5/3-VL, Qwen3.5,\\ Florence-2, InternVL, Gemma\\ 或 Cosmos / Wan（视频 DiT）};
    \node[box, text width=1.2cm, right=0.5cm of vlm, fill=cbPurple!15] (h) {hidden\\ state $z$};
    \node[box, text width=2.0cm, right=0.5cm of h, fill=emphasis!12] (head) {动作头\\ \scriptsize FAST / OFT / $\pi$ / GR00T};
    \node[box, text width=1.5cm, right=0.5cm of head, fill=neutral!15] (act) {动作块\\ \scriptsize 归一化动作};
    \draw[arr] (obs) -- (vlm); \draw[arr] (vlm) -- (h); \draw[arr] (h) -- (head); \draw[arr] (head) -- (act);
    \draw[decorate, decoration={brace, amplitude=5pt, mirror, raise=3pt}] (obs.south west) -- (act.south east)
      node[midway, below=9pt, font=\scriptsize] {外层契约：\texttt{forward} / \texttt{predict\_action} 都直接吃原始观测，训练输入 = 部署观测};
    \draw[decorate, decoration={brace, amplitude=5pt, raise=3pt}] (vlm.north west) -- (head.north east)
      node[midway, above=9pt, font=\scriptsize] {内层契约：骨干与头通过统一 hidden-state 规格连接，可独立替换};
  \end{tikzpicture}
  \end{center}
  \vspace{4pt}
  \begin{itemize}
    \item 每个 framework 文件可单独运行（smoke test）；YAML + 点路径命令行覆盖（\texttt{--framework.qwenvl.base\_vlm}）
    \item 冻结 = 模块名正则；分模块学习率 = \texttt{trainer.learning\_rate} 字典
    \item 评测：WebSocket 策略服务器 + 基准侧薄适配器（\texttt{model2libero\_interface.py} 等），仿真与真机共用
  \end{itemize}
\end{frame}

% ---------------------------------------------------------------- 5
\begin{frame}{StarVLA：四种动作头，只在"如何从 $z$ 取动作"上不同}
  \vspace{2pt}
  \begin{center}\small
  \begin{tabular}{llll}
    \toprule
    头 & 接口 & 目标 & 推理 \\
    \midrule
    FAST & 离散 token，LLM 词表 & $-\sum_m \log p_\theta(z_m \mid H_t, z_{<m})$ & 自回归 + 逆 tokenizer \\
    OFT & $K$ 个 action query & $\frac{1}{Kd_a}\sum_k \| g_\phi(h^{\mathrm{act}}_{t,k}) - a_{t+k}\|_1$ & 一次前向 \\
    $\pi$ & 逐层 cross-DiT 流匹配 & $\mathbb{E}\|v_\phi(A^\tau,\tau;H_t)-(A-\epsilon)\|_2^2$ & $N$ 步积分 \\
    GR00T & 双系统：VLM 慢 + DiT 快 & 同上，额外条件于状态 $s_t$、本体 $e$ & $N$ 步去噪 \\
    \bottomrule
  \end{tabular}
  \end{center}
  \vspace{4pt}
  线性路径 $A^\tau=(1-\tau)\epsilon+\tau A$，目标速度 $u=A-\epsilon$；$\pi$ 与 GR00T 共享这一目标，差别是运动模块的耦合方式。
  \vspace{4pt}
  \begin{itemize}
    \item \pos{OFT} 同时是诊断工具：表现好 $\Rightarrow$ 骨干表征线性可读地暴露了连续控制信息
    \item \con{FAST} 是唯一不改架构的方案，代价是离散瓶颈 + 串行解码
    \item 新范式 = 实现并注册一个头；骨干、训练循环、评测管线不动
  \end{itemize}
\end{frame}

% ---------------------------------------------------------------- 6
\begin{frame}{StarVLA：训练模式、基准生态与扩展效率}
  \begin{columns}[T]
    \column{0.50\textwidth}
    \textbf{三种训练模式 = 三个脚本}
    \begin{center}\footnotesize
    \begin{tabular}{p{3.0cm}p{3.9cm}}
      \toprule
      模式 & 入口 \\
      \midrule
      行为克隆 SFT & \texttt{train\_starvla.py} \\
      多模态共训（双 loader） & \texttt{train\_starvla\_cotrain.py} \\
      跨本体混合 & \texttt{datasets.vla\_data.data\_mix} \\
      VLM 单独训练 & \texttt{train\_starvlm.py} \\
      \bottomrule
    \end{tabular}
    \end{center}
    \vspace{4pt}
    \textbf{13 个仿真基准} LIBERO / LIBERO-plus / RoboTwin 2.0 / RoboDojo / VLA-Arena / DOMINO / RoboCasa (GR1, 365) / SimplerEnv / BEHAVIOR / CALVIN / MetaWorld / VLN-CE；真机 Franka / Realman / Unitree G1 / RoboChallenge
    \column{0.46\textwidth}
    \textbf{多节点扩展效率}（GR00T 头，每 GPU batch 8，A100）
    \begin{center}\footnotesize
    \begin{tabular}{rrrr}
      \toprule
      GPU & 秒/步 & samples/s & 效率 \\
      \midrule
      8 & 0.735 & 87 & 100\% \\
      32 & 0.899 & 285 & 81.9\% \\
      64 & 0.925 & 554 & 79.6\% \\
      256 & 0.931 & 2200 & 79.1\% \\
      \bottomrule
    \end{tabular}
    \end{center}
    \vspace{2pt}
    \footnotesize 跨节点一次性 +0.2 s/步，之后平台化；固定步数不会因加卡变快，数据量驱动的训练可放心扩到数百卡。
  \end{columns}
\end{frame}

% ---------------------------------------------------------------- 7
\section{StarVLA-$\alpha$}
\begin{frame}{StarVLA-$\alpha$：最简基线 = Qwen3-VL-4B + MLP 头}
  \textbf{设计：}原始 RGB + 指令；训练集 z-score 归一化；无状态、无历史、无机器人预训练；骨干 lr 1e-5、头 lr 1e-4、cosine、$\le$100k 步。
  \vspace{4pt}
  \begin{center}\footnotesize
  \begin{tabular}{lrrrrr}
    \toprule
    方法 & LIBERO avg & WidowX & Google VM & RoboTwin clean$^*$ / random$^*$ & RoboCasa-GR1 \\
    \midrule
    OpenVLA-OFT & 97.1 & 31.3 & 63.0 & – & – \\
    $\pi_0$ & 94.1 & 27.1 & 58.8 & 65.9 / 58.4 & – \\
    $\pi_{0.5}$ & 96.9 & 46.9 & 72.7 & 82.7 / 76.8 & 37.0 \\
    GR00T-N1.6 & 97.0 & 62.0 & 67.7 & – & 47.6 \\
    \textbf{StarVLA-$\alpha$}（specialist） & \HL{98.8} & 64.6 & \HL{76.0} & 88.2 / 88.3 & 53.8 \\
    \textbf{StarVLA-$\alpha$}（generalist） & 97.8 & \HL{65.2} & 74.3 & \HL{88.7} / 87.8 & \HL{57.3} \\
    \bottomrule
  \end{tabular}
  \end{center}
  \vspace{2pt}
  {\footnotesize $^*$ clean + random 数据都用于训练。}
  \begin{itemize}
    \item 唯一输掉的一格：RoboTwin Base（仅 50 clean/任务）50.3 $<$ $\pi_{0.5}$ 60.2 —— \con{低数据 + 关节角控制是短板}
    \item 四个基准数据合训一个 generalist，GR1 反而 +3.5
  \end{itemize}
\end{frame}

% ---------------------------------------------------------------- 8
\begin{frame}{StarVLA-$\alpha$：常识一，动作头设计重要吗？}
  同一骨干、同一数据、同一训练协议，只换头：
  \vspace{4pt}
  \begin{center}\small
  \begin{tabular}{lrrrr}
    \toprule
    头 & LIBERO avg & WidowX & RoboTwin clean$^*$ & RoboCasa-GR1 \\
    \midrule
    MLP（OFT 式） & 98.8 & 64.6 & 88.2 & \HL{53.8} \\
    FAST（离散 token） & 97.8 & \con{35.6} & \con{72.5} & 45.0 \\
    GR00T（双系统流匹配） & 98.7 & 65.3 & 88.0 & 52.8 \\
    $\pi$（流匹配专家） & 98.1 & 65.9 & 88.1 & 48.9 \\
    \bottomrule
  \end{tabular}
  \end{center}
  \vspace{6pt}
  \begin{itemize}
    \item \textbf{连续 $\gg$ 离散}：FAST 在 WidowX 落后 29 个点、RoboTwin 落后 16 个点
    \item \textbf{三种连续头相当}：LIBERO 与 RoboTwin 差距 $<$ 1 个点
    \item 结论：骨干够强时，动作头的额外复杂度不必要 —— \mute{但这是"同头微调"的结论，VLAct 会追问"骨干可复用吗"}
  \end{itemize}
\end{frame}

% ---------------------------------------------------------------- 9
\begin{frame}{StarVLA-$\alpha$：常识二，现有动作预训练重要吗？}
  \begin{center}\small
  \begin{tabular}{lrrrrr}
    \toprule
    中间预训练 & 轨迹数 & RoboTwin 50$\times$50 & +Random$\times$500 & GR1 24$\times$10 & GR1 24$\times$1000 \\
    \midrule
    无 & – & 50.3 & 88.2 & 9.8 & 53.8 \\
    + OXE（跨域） & 232.6k & \con{30.2} & 83.6 & \con{1.2} & \con{27.8} \\
    + InternData-A1（同本体） & 630k & \HL{63.6} & 88.6 & 2.8 & 35.4 \\
    + RoboTwin-Rand（同域） & 25k & \HL{79.7} & 88.8 & 2.2 & 33.3 \\
    \bottomrule
  \end{tabular}
  \end{center}
  \vspace{6pt}
  \begin{alertblock}{双刃剑}
    跨域数据全面拖后腿；同域数据只在目标域低数据段有用，换到 GR1 一律掉点。作者建议"谨慎使用"。
  \end{alertblock}
  \vspace{4pt}
  这张表是 VLAct 问题定义的直接来源：\textbf{是数据不对，还是训练方式不对？}
\end{frame}

% ---------------------------------------------------------------- 10
\begin{frame}{StarVLA-$\alpha$：常识三 + generalist 的三个发现}
  \begin{columns}[T]
    \column{0.48\textwidth}
    \textbf{数据工程}（本体感受 / 历史帧 / delta / relative）
    \begin{itemize}
      \item 低数据段 +1$\sim$10 点（proprio 在 RoboTwin 50$\times$50 +10.5）
      \item 数据充足后全部归零；\con{历史帧多数设定略降}
    \end{itemize}
    \vspace{4pt}
    \textbf{多本体动作参数化}（GR1 成功率）
    \begin{center}\footnotesize
    \begin{tabular}{lr}
      \toprule
      RDT 物理统一空间 & 52.3 \\
      每本体独立头 & 53.5 \\
      \textbf{零填充到 32 维} & \HL{57.3} \\
      \bottomrule
    \end{tabular}
    \end{center}
    \column{0.48\textwidth}
    \textbf{规模与优化}（generalist 设定）
    \begin{center}\footnotesize
    \begin{tabular}{p{1.2cm}p{2.7cm}p{2.3cm}}
      \toprule
      变量 & 设定 & GR1 \\
      \midrule
      初始化 & 随机 / Florence-2 / Qwen3-VL & 28.8 / 39.2 / 57.3 \\
      模型规模 & 2B / 4B / 8B & 50.7 / 57.3 / 58.2 \\
      batch & 64 / 256 / 1024 & 40.0 / 48.3 / \HL{59.2} \\
      \bottomrule
    \end{tabular}
    \end{center}
    \begin{itemize}
      \item 4B$\to$8B 增益 $<$1 点，\pos{4B 是甜点}
      \item batch 64$\to$1024 使 GR1 +19 点，\textbf{比模型规模影响更大}
      \item 真机 RoboChallenge ARX5：SR 33.6 vs $\pi_{0.5}$ 12.7
    \end{itemize}
  \end{columns}
\end{frame}

% ---------------------------------------------------------------- 11
\section{VLAct}
\begin{frame}{VLAct：先导实验发现 decoder lock-in}
  固定 Qwen3-VL-4B，只改预训练头与微调头（LIBERO-Plus / RoboTwin-Clean）：
  \vspace{4pt}
  \begin{columns}[T]
    \column{0.50\textwidth}
    \textbf{(a) 离散监督能迁移，但丢信息}
    \begin{itemize}
      \item FAST 预训练 $\to$ GR00T 微调：略好于从零
      \item 保留 FAST 头：远低于任何连续头，预训练补不回来
      \item 离散化丢掉了幅度与时序细节
    \end{itemize}
    \vspace{4pt}
    \textbf{(b) 单一连续头造成 head-specific 坍缩}
    \begin{itemize}
      \item OFT 预训练 $\to$ OFT 微调：大涨
      \item 同一骨干 $\to$ PI / GR00T 微调：\con{低于从零}
      \item 动作信息还在，但只有原头能解码
    \end{itemize}
    \column{0.46\textwidth}
    \begin{center}\small
    \begin{tabular}{lrr}
      \toprule
      预训练头 & 下游 PI 微调 & vs 从零 \\
      \midrule
      无 & 60.5 & – \\
      OFT & \con{55.1} & $-5.4$ \\
      OFT + GR00T & 63.1 & $+2.6$ \\
      OFT + PI + GR00T & \HL{77.0} & $+16.5$ \\
      \bottomrule
    \end{tabular}
    \end{center}
    \vspace{6pt}
    \begin{exampleblock}{命名}
      \textbf{decoder lock-in}：骨干把动作信息组织成预训练那个头的解码几何。同头成功率高估了骨干可复用性。
    \end{exampleblock}
  \end{columns}
\end{frame}

% ---------------------------------------------------------------- 12
\begin{frame}{VLAct：配方总览，三个组件只在持续预训练阶段使用}
  \begin{center}
  \begin{tikzpicture}[node distance=0.45cm]
    \node[box, text width=2.7cm] (c1) {\textbf{① 保护 VLM 先验}\\ \scriptsize 冻视觉编码器 + LLM 下半层\\ caption 混训 $0.5\,\mathcal{L}_{\text{VLM-CE}}$};
    \node[box, text width=2.7cm, right=0.5cm of c1] (c2) {\textbf{② 多头共监督}\\ \scriptsize OFT + PI + GR00T\\ 同一 $z$、同一目标 $a$};
    \node[box, text width=2.7cm, right=0.5cm of c2] (c3) {\textbf{③ 部分统一动作空间}\\ \scriptsize 20 维布局，夹爪共享\\ 非激活维 mask，wrap loss};
    \node[box, text width=3.2cm, right=0.8cm of c3, fill=emphasis!12] (ft) {\textbf{④ 下游微调}\\ \scriptsize 丢弃预训练头与 caption 流\\ 新初始化任务头，全参数解冻\\ 数据 / 优化器 / 预算与基线相同};
    \draw[arr] (c3) -- node[above, font=\scriptsize] {只交付骨干} (ft);
    \node[draw, dashed, rounded corners, fit=(c1)(c2)(c3), inner sep=6pt, label={[font=\scriptsize]above:持续预训练（DROID + InternData-A1 + RoboCoin + MolmoAct + caption，16 GPU）}] {};
  \end{tikzpicture}
  \end{center}
  \vspace{4pt}
  \[
    \mathcal{L}_{\text{total}} = \underbrace{\mathcal{L}_{\text{OFT}} + \mathcal{L}_{\text{PI}} + \mathcal{L}_{\text{GR00T}}}_{\mathcal{L}_{\text{action}}\ (\text{每头再加 } \mathcal{L}_{\text{wrap}})} + 0.5\,\mathcal{L}_{\text{VLM-CE}}
  \]
  \vspace{2pt}
  \begin{itemize}
    \item 所有对比中\textbf{唯一变量是骨干权重}：7.6$\sim$21.4 个点的增益归因到配方
    \item 无新模块、无本体适配器、无路由：可在 StarVLA 上直接复现
  \end{itemize}
\end{frame}

% ---------------------------------------------------------------- 13
\begin{frame}{VLAct 组件 1--2：先验保护与多头共监督的消融}
  \begin{columns}[T]
    \column{0.48\textwidth}
    \textbf{① 冻结策略}（预训练阶段）
    \begin{center}\footnotesize
    \begin{tabular}{lrr}
      \toprule
      更新范围 & LIBERO-Plus & RoboTwin \\
      \midrule
      全骨干 & 78.9 & 77.1 \\
      冻视觉编码器 & 81.3 & 79.3 \\
      冻视觉编码器 + 下半 LLM & \HL{82.6} & \HL{80.5} \\
      \bottomrule
    \end{tabular}
    \end{center}
    \vspace{4pt}
    \textbf{① 辅助数据}（LIBERO-Plus）
    \begin{itemize}
      \item robot-only 75.0 $\to$ +caption \HL{82.6}
      \item Spatial-QA / Point-QA / BBox-QA 均有提升
      \item \pos{纯文本指令数据也提升}：增益只能解释为"保住基础模型工作区间"
      \item 全混合 82.5：固定步数下稀释了 caption
    \end{itemize}
    \column{0.48\textwidth}
    \textbf{② 多头 vs 单头，同头微调}（RoboTwin Base Clean）
    \begin{center}\footnotesize
    \begin{tabular}{lrrrr}
      \toprule
      下游头 & 从零 & 单头 & 三头 & $\Delta$ \\
      \midrule
      OFT & 61.7 & 78.8 & \HL{80.5} & +1.7 \\
      PI & 60.5 & 75.4 & \HL{77.0} & +1.6 \\
      GR00T & 51.2 & 71.7 & \HL{76.0} & +4.3 \\
      \bottomrule
    \end{tabular}
    \end{center}
    \vspace{6pt}
    \begin{itemize}
      \item 头多样性 = 正则化：同头也涨，GR00T 受益最大
      \item 三头共享一次骨干前向，只多头本身的计算 \mute{（论文未给显存 / 吞吐数字）}
    \end{itemize}
  \end{columns}
\end{frame}

% ---------------------------------------------------------------- 14
\begin{frame}{VLAct 组件 3：部分统一动作空间与 wrap-aware loss}
  \begin{columns}[T]
    \column{0.50\textwidth}
    \textbf{20 维固定输出，按本体角色分配}
    \begin{center}\footnotesize
    \begin{tabular}{p{0.9cm}p{3.7cm}p{1.9cm}}
      \toprule
      维度 & 含义 & 表示 \\
      \midrule
      1–6 & 双臂左臂（AgileX） & 绝对关节角 \\
      7–12 & 双臂右臂 & 绝对关节角 \\
      13–18 & 单臂（Franka） & delta 末端位姿 \\
      \textbf{19} & \textbf{共享夹爪}：Franka 单夹爪 = AgileX 左夹爪 & [0,1] \\
      20 & 双臂右夹爪 & [0,1] \\
      \bottomrule
    \end{tabular}
    \end{center}
    每个样本只在激活维上算 loss，其余 mask。
    \vspace{4pt}
    \begin{center}\footnotesize
    \begin{tabular}{lrr}
      \toprule
      动作空间 & RoboTwin & LIBERO-Plus \\
      \midrule
      每本体独立头 & 78.5 & 81.1 \\
      朴素统一头 & 79.5 & 81.4 \\
      \textbf{部分统一} & \HL{80.5} & \HL{82.6} \\
      \bottomrule
    \end{tabular}
    \end{center}
    \column{0.46\textwidth}
    \textbf{周期关节角}：179° 与 $-179^\circ$ 物理上只差 2°
    \[
      a_{\text{wrap}} = (a+\pi) \bmod 2\pi - \pi
    \]
    \[
      \delta_{\text{wrap}} = \big((\hat a - a)+\pi\big) \bmod 2\pi - \pi,\quad
      \mathcal{L}_{\text{wrap}} = |\delta_{\text{wrap}}|
    \]
    加到每个头的目标上；对 PI / GR00T 作用于最终生成样本；只用于绝对关节角维度。
    \vspace{4pt}
    \begin{center}\footnotesize
    \begin{tabular}{lr}
      \toprule
      设定（RoboTwin Base Clean） & 成功率 \\
      \midrule
      原始回归 & 75.5 \\
      数据侧 wrap & 78.6 \\
      + wrap-aware loss & \HL{80.5} \\
      \bottomrule
    \end{tabular}
    \end{center}
    \pos{+5 个点，零成本}，VLAct 单项消融里最大的一项。
  \end{columns}
\end{frame}

% ---------------------------------------------------------------- 15
\begin{frame}{VLAct：主结果，只换骨干权重}
  \begin{center}\footnotesize\setlength{\tabcolsep}{4pt}
  \begin{tabular}{llrrl}
    \toprule
    基准 & 设定 & VLAct & 同骨干基线 & 最强外部对手 \\
    \midrule
    LIBERO-Plus & 7 类扰动均值 & \HL{82.6} & 75.0 & ABot-M0 80.5 \\
    VLA-Arena & 11 套件加权 & \HL{54.8} & 33.4 & $\pi_{0.5}$ 44.3 \\
    RoboTwin 2.0 Base & Clean / Random & \HL{80.5 / 41.5} & 61.7 / 10.5 & X-VLA 70.0 / 39.0 \\
    RoboTwin 2.0 Data Scaling & Clean / Random & \HL{92.5} / 90.8 & 88.2 / 88.3 & HoloBrain-0 91.9 / 92.3 \\
    DOMINO（动态） & SR / MS & \HL{18.50 / 34.20} & 10.86 / 30.49 & $\pi_{0.5}$ 9.63 / 26.17 \\
    RoboCasa-GR1（未见本体） & 全量 & \HL{54.0} & 48.8 & GR00T-N1.6 47.6 \\
    RoboCasa-GR1 & \textbf{20\% 数据} & \HL{49.5} & – & 超全量 GR00T-N1.6 \\
    RoboDojo（ARX X5，未见） & 分数 / 成功率 & 10.66 / 7.60 & $\alpha$ 6.40 / 3.24 & DM0.5 24.90 / 19.34 \\
    真机 Franka 单臂短程 & 4 任务 & \HL{92.5} & 77.5 & – \\
    真机 Franka 双臂协调 & 5 任务 & \HL{72.0} & 44.0 & – \\
    \bottomrule
  \end{tabular}
  \end{center}
  \vspace{4pt}
  \begin{itemize}
    \item LIBERO-Plus 增益集中在 Camera（+26.9）/ Robot / Noise / Layout；\con{Language 维度 $-5.5$}
    \item 预训练只用单臂数据，双臂协调 44.0 $\to$ 72.0（叠裤子、叠毛巾各 +30）
    \item RoboDojo 35 策略中成功率第 6，超过全部 4 个 WAM；\con{Memory 0.66 / 0.56 接近零}
  \end{itemize}
\end{frame}

% ---------------------------------------------------------------- 16
\begin{frame}{对话：双刃剑 vs 配方}
  \begin{center}\footnotesize
  \begin{tabular}{p{2.3cm}p{4.6cm}p{6.4cm}}
    \toprule
     & StarVLA-$\alpha$（2026-04） & VLAct（2026-08） \\
    \midrule
    骨干 & Qwen3-VL-4B & Qwen3-VL-4B \\
    预训练数据 & OXE / InternData-A1 / RoboTwin-Rand & DROID + InternData-A1 + RoboCoin + MolmoAct + caption \\
    预训练方式 & 单 MLP 头，全参数，动作-only & 冻浅层，三头，部分统一空间，caption 混训 \\
    GR1（未见本体） & 53.8 $\to$ \con{35.4}（+A1） & 48.8 $\to$ \HL{54.0}；20\% 数据 49.5 \\
    RoboTwin Base Clean & 50.3 $\to$ 63.6（+A1） & 61.7 $\to$ 80.5 \\
    结论 & 预训练是双刃剑 & 预训练是否有用取决于配方 \\
    \bottomrule
  \end{tabular}
  \end{center}
  \vspace{6pt}
  \begin{itemize}
    \item "Beyond Data Scaling" 的实证含义：\textbf{同样的数据源，配方决定符号}
    \item 动作空间上两篇也叠加：零填充 $>$ 独立头 / RDT（$\alpha$ 表 6）；零填充之上再做物理同义维对齐再 +1（VLAct 表 12）
    \item 共训保护先验一线：ST4VLA（StarVLA 报告 §6）Google VM 66.1 $\to$ 84.6，VLAct 选 caption 作最强单源锚点
  \end{itemize}
\end{frame}

% ---------------------------------------------------------------- 17
\begin{frame}{批判性评价}
  \begin{columns}[T]
    \column{0.48\textwidth}
    \textbf{\pos{做得好}}
    \begin{itemize}
      \item 归因干净：头、初始化、数据、优化器、预算全同，只换骨干
      \item 组件简单到可以被推翻：一个正则冻结、三个已有头并联、一张 20 维表
      \item 把"哪个头最好"换成"如何让骨干对头不敏感"
      \item 成本诚实：开源数据 + 16 GPU，与工业系统差距明标
    \end{itemize}
    \column{0.48\textwidth}
    \textbf{\con{局限}}
    \begin{itemize}
      \item 只有 4B；持续预训练下的规模曲线空白
      \item Memory 接近零：单帧、无历史的架构假设没被挑战
      \item "头无关"有边界：Base/Random 下 OFT 41.5 vs GR00T 22.9 / PI 23.7
      \item Language 维度掉 5.5 点
      \item Data Scaling 对比算力不可控；真机每任务 10 次 rollout
      \item 三头共训开销未量化；20 维布局手工设计，人形 / 灵巧手如何扩展未讨论
      \item 同名基线不一致：RoboTwin Base OFT 61.7（VLAct）vs 50.38（StarVLA README）
    \end{itemize}
  \end{columns}
\end{frame}

% ---------------------------------------------------------------- 18
\section{代码库}
\begin{frame}{StarVLA 代码库地图（\texttt{starVLA\_dev} @ d81fc66，2026-09）}
  \begin{columns}[T]
    \column{0.52\textwidth}
    \begin{center}\footnotesize
    \begin{tabular}{p{3.4cm}p{3.6cm}}
      \toprule
      目录 & 职责 / 关键类 \\
      \midrule
      \texttt{model/framework/} & 28 个注册框架（VLM4A 20、WM4A 6、VM4A 2）；\texttt{build\_framework(cfg)} 查注册表 \\
      \texttt{modules/action\_model/} & 11 个动作头文件：MLP(OFT)、FAST、Layerwise-FM(PI)、FM DiT(GR00T) \\
      \texttt{modules/vlm/} & 10 个接口，\texttt{get\_vlm\_model} 分派 9 个分支 \\
      \texttt{dataloader/} & GR00T LeRobot 管线移植；VLM LLaVA-json；UMI 适配 \\
      \texttt{training/} & 4 个显式循环脚本 + \texttt{trainer\_utils}（冻结 / 分组 lr / 加载） \\
      \texttt{examples/} & 25 个目录：13 仿真基准、5 真机、6 模型扩展、1 human2robot \\
      \texttt{deployment/} & WebSocket / ZMQ 策略服务器、Docker、HF 上传 \\
      \bottomrule
    \end{tabular}
    \end{center}
    \column{0.44\textwidth}
    \textbf{实测规模}
    \begin{itemize}
      \item 261 个 Python 文件，57.6k 行；\texttt{starVLA/} 139 个文件
      \item 12 个 \texttt{model2*\_interface.py} 基准适配器
      \item 测试：7 文件 / 40 用例 / 1,067 行，\con{无 CI}
      \item 浅历史 131 次提交，40 位作者
    \end{itemize}
    \vspace{4pt}
    \textbf{契约在代码中的落点}
    \begin{itemize}
      \item \texttt{baseframework}（继承 \texttt{PreTrainedModel}）软约束 \texttt{forward} / \texttt{predict\_action}
      \item dataloader 的 \texttt{collate\_fn} 就是 \texttt{return batch}：样本 = \{action, image, lang, robot\_tag, state?\}
      \item 每个框架文件末尾带 \texttt{\_\_main\_\_} smoke test
      \item 三份训练循环是复制粘贴 —— \con{最明显的可重构点}
    \end{itemize}
  \end{columns}
\end{frame}

% ---------------------------------------------------------------- 19
\begin{frame}{VLAct 配方在 StarVLA 代码中的状态}
  \begin{center}\footnotesize\setlength{\tabcolsep}{4pt}
  \begin{tabular}{p{0.4cm}p{3.3cm}p{0.9cm}p{4.2cm}p{3.8cm}}
    \toprule
    \# & 配方 & 状态 & 代码现状 & 需改动 \\
    \midrule
    a & 冻视觉编码器 + LLM 下半层 & \con{部分} & \texttt{trainer.freeze\_modules} 只接受精确点路径；4B 有 36 层 & 列 18 条路径即可；建议加正则 / 区间语法 \\
    b & caption 共训 $0.5\,\mathcal{L}_{\text{VLM-CE}}$ & \HL{已有} & \texttt{train\_starvla\_cotrain.py} + \texttt{loss\_scale.vlm}；已注册 \texttt{sharegpt4v\_coco} & 改 yaml；VLAct 的 caption 源需注册 \\
    c & OFT + PI + GR00T 三头共监督 & \con{缺失} & 28 个框架全部单头；训练器只读 \texttt{action\_loss} & 新框架 \texttt{QwenMultiHead}：一次骨干前向，三头 loss 相加 \\
    d & 20 维部分统一布局 + mask & \con{部分} & mask 消费端：\texttt{masked\_l1\_loss}；生产端：UMI / MiniCPM 80 维；FM 头无 mask & 新 transform + 两个 FM 头加 mask \\
    e & 周期关节 wrap-aware L1 & \con{缺失} & 无任何模 $2\pi$ 逻辑 & 新 loss 函数 + 数据侧 wrap \\
    f & 丢预训练头、重初始化、全参解冻 & \HL{已有} & \texttt{pretrained\_checkpoint} + \texttt{reload\_modules} + \texttt{freeze\_modules: ''} & 注意 PI 的 \texttt{project\_layers} \\
    \bottomrule
  \end{tabular}
  \end{center}
  \vspace{4pt}
  \begin{itemize}
    \item 两项"已有"、两项"部分"、两项"缺失"：复现 VLAct 的工程量集中在 \textbf{多头框架} 与 \textbf{动作空间 transform}
    \item 唯一的多 loss 先例：LangForce 在同一 GR00T 头上加权两支路 loss
    \item 完整的 diff 级改动建议见 \texttt{reports/02\_starvla\_codebase\_analysis.md} §9
  \end{itemize}
\end{frame}

% ---------------------------------------------------------------- 20
\section{基准与路线图}
\begin{frame}{基准选择：研究"持续预训练 / 表示学习"该用哪几个}
  \begin{center}\footnotesize\setlength{\tabcolsep}{4pt}
  \begin{tabular}{p{0.9cm}p{2.7cm}p{3.3cm}p{6.6cm}}
    \toprule
    优先级 & 基准 & 考察的能力 & 理由 \\
    \midrule
    1 & LIBERO-plus & 零样本鲁棒性（7 维扰动） & 训练集固定，提升只能来自骨干；10,030 实例噪声小 \\
    2 & RoboTwin 2.0 \textbf{Base} & 少样本 + clean$\to$random & 50 条/任务放大表征价值；必须同时报 clean 与 random \\
    3 & RoboCasa-GR1 数据比例曲线 & 未见本体样本效率 & 人形，形态 / 动作空间都不同；基线也要跑同样比例点 \\
    4 & VLA-Arena & L0$\to$L1/L2 外推 + 安全代价 & 与"表示能否支撑结构外推"一致 \\
    5 & RoboDojo & 第三方裁判 + 五维诊断 & 唯一带 Memory / Precision 拆分的榜单；适合最终报告 \\
    \midrule
    \con{不建议} & LIBERO 标准版、SimplerEnv、MetaWorld、CALVIN D$\to$D、BEHAVIOR & – & 饱和、方差大、或评测成本极高 \\
    \bottomrule
  \end{tabular}
  \end{center}
  \vspace{4pt}
  \textbf{报告协议}：固定下游预算并写明 checkpoint 规则；$\ge$3 seeds；区分见过 / held-out 本体（VLAct 只有 GR1 与 ARX X5 是真 held-out）；RoboTwin 四个数都给；附表示层诊断（跨头迁移矩阵）。
  \vspace{4pt}
  \textbf{没有好基准的维度}：Memory（RoboDojo 6 任务几乎人人接近零）、中等长程分步计分、跨本体零样本、动态场景（只有 DOMINO）。
\end{frame}

% ---------------------------------------------------------------- 21
\begin{frame}{研究路线图：六类 18 个方向}
  \begin{center}\footnotesize
  \begin{tabular}{p{2.2cm}p{5.4cm}p{6.0cm}}
    \toprule
    类别 & 方向 & 出发点的证据 \\
    \midrule
    A 表征诊断 & A1 骨干可复用性探针套件；A2 自动化"哪层该冻" & decoder lock-in 只能间接观察；冻结只有 3 档消融 \\
    B 预训练目标 & B1 头多样性推广（FAST / 未来帧 / 空间 QA 头）；B2 辅助数据调度；B3 潜动作 vs 多头；B4 世界模型辅助头 & 三头同头 +1.6$\sim$4.3；全混合被稀释；"广义 VLA 视角" \\
    C 动作空间 & C1 可学习的部分统一布局；C2 几何一致参数化（SO(3)）；C3 零样本本体迁移 & 20 维手工表；wrap loss +5；GR1 20\% 才 49.5 \\
    D 能力短板 & D1 记忆与长程；D2 语言泛化；D3 低数据下生成式头落后 & Memory 0.66；Language $-5.5$；Random 下 OFT 41.5 vs 22.9 \\
    E 系统 scaling & E1 规模曲线；E2 三头开销与降本；E3 数据 $\times$ 配方 Pareto & 只有 4B；开销未量化 \\
    F 评测方法 & F1 骨干基准协议；F2 generalist $\times$ 持续预训练；F3 真机统计功效 & 只换骨干未标准化；两条线未合并；n=10 \\
    \bottomrule
  \end{tabular}
  \end{center}
  \vspace{4pt}
  详细的做法、代码落点、评测方案与风险见 \texttt{reports/07\_research\_roadmap.md}。
\end{frame}

% ---------------------------------------------------------------- 22
\begin{frame}{优先级与六个月计划（16 GPU）}
  \begin{columns}[T]
    \column{0.42\textwidth}
    \textbf{影响 $\times$ 成本}
    \begin{center}\footnotesize
    \begin{tabular}{lll}
      \toprule
      方向 & 影响 & 成本 \\
      \midrule
      \HL{A1 诊断套件} & 高 & 低 \\
      \HL{F1 骨干基准} & 高 & 低 \\
      \HL{E2 三头开销} & 中 & 低 \\
      B1 头多样性 & 高 & 中 \\
      D1 记忆 & 高 & 中 \\
      C1 可学习动作空间 & 高 & 中 \\
      B4 世界模型辅助头 & 高 & 高 \\
      E1 规模曲线 & 高 & 很高 \\
      \bottomrule
    \end{tabular}
    \end{center}
    \column{0.56\textwidth}
    \textbf{执行计划}
    \begin{center}\footnotesize
    \begin{tabular}{cp{6.2cm}}
      \toprule
      月 & 目标与交付 \\
      \midrule
      1 & 复现 VLAct；探针 + 骨干基准脚本；三头开销表 \\
      2 & A2 自动冻结；B1 头组合因子实验 $\to$ 跨头迁移矩阵 \\
      3 & D1 记忆三路线小规模验证；D3 步数曲线 \\
      4 & C1 可学习布局，加 GR1 做三本体预训练 \\
      5 & B2 调度器；B4 特征预测辅助头（DOMINO / Random） \\
      6 & 整合最优组合，全基准评测 + RoboDojo 提交，成文并向 StarVLA 提 PR \\
      \bottomrule
    \end{tabular}
    \end{center}
    \footnotesize 每月实验都遵守 F1 协议（只换骨干）；同时报告"同头"与"跨头"两组数字。
  \end{columns}
\end{frame}

% ---------------------------------------------------------------- 23
\begin{frame}{三个 takeaway}
  \begin{enumerate}
    \item \textbf{VLM 骨干是 VLA 的一阶设计变量。} 只换骨干权重就能拿到 7.6$\sim$21.4 个点；朴素动作拟合会把它变成负资产（GR1 9.8 $\to$ 1.2），保护先验 + 多头 + 统一动作空间把它变成净收益（20\% 数据超全量 GR00T-N1.6）。
    \vspace{6pt}
    \item \textbf{头多样性是一种廉价的正则化。} 单头预训练造成 decoder lock-in（PI 60.5 $\to$ 55.1），加一个头就翻正，三头对同头也涨 1.6$\sim$4.3。这个原则可以推到非动作头。
    \vspace{6pt}
    \item \textbf{StarVLA 让"只改一个变量"成为默认工作方式。} 两条契约 + 13 个基准 + 79\% 的 256 卡扩展效率；下一步最值钱的是诊断工具、骨干基准协议，以及所有人都接近零的 Memory 维度。
  \end{enumerate}
  \vspace{8pt}
  \begin{center}
    \small 回到开场的问题：预训练不是双刃剑，\textbf{配方决定符号}。
  \end{center}
\end{frame}

% ---------------------------------------------------------------- 24
\begin{frame}{References}
  \small
  \begin{thebibliography}{9}
    \bibitem{vlact} S. Yang, C. Wang, et al. \emph{Beyond Data Scaling: Representation-Centric Continued Pre-training for Vision-Language-Action Models.} arXiv:2608.27550, 2026.
    \bibitem{alpha} J. Ye, N. Gao, S. Yang, et al. \emph{StarVLA-$\alpha$: Reducing Complexity in Vision-Language-Action Systems.} arXiv:2604.11757, ECCV 2026.
    \bibitem{starvla} StarVLA Community. \emph{StarVLA: A Lego-like Codebase for Vision-Language-Action Model Developing.} arXiv:2604.05014, 2026. github.com/starVLA/starVLA
    \bibitem{oft} M. J. Kim, C. Finn, P. Liang. \emph{Fine-Tuning Vision-Language-Action Models: Optimizing Speed and Success (OpenVLA-OFT).} arXiv:2502.19645, 2025.
    \bibitem{pi0} K. Black et al. \emph{$\pi_0$: A Vision-Language-Action Flow Model for General Robot Control.} arXiv:2410.24164, 2024.
    \bibitem{groot} J. Bjorck et al. \emph{GR00T N1: An Open Foundation Model for Generalist Humanoid Robots.} arXiv:2503.14734, 2025.
    \bibitem{liberoplus} S. Fei et al. \emph{LIBERO-Plus: In-depth Robustness Analysis of Vision-Language-Action Models.} arXiv:2510.13626, 2025.
    \bibitem{robodojo} T. Chen et al. \emph{RoboDojo: A Unified Sim-and-Real Benchmark for Generalist Robot Manipulation Policies.} arXiv:2607.04434, 2026.
  \end{thebibliography}
\end{frame}

% ---------------------------------------------------------------- 25
\begin{frame}
  \begin{center}
    {\Large 谢谢}\\[10pt]
    \normalsize 仓库：\url{https://github.com/asimfish/awesome_starvla}\\[4pt]
    \small 3 篇论文中英 PDF · 7 份精读报告 · 文献编目 · 本 PPT 源码
  \end{center}
\end{frame}

% ================================================================ backup
\appendix
\begin{frame}{Backup Slides}
  \begin{center}\Large 备份页\end{center}
\end{frame}

\begin{frame}{Backup：VLAct 数据清洗规则（附录 H）}
  \begin{itemize}
    \item 本体与动作空间：Franka 7 维（6 delta EE + 1 夹爪）；AgileX 14 维（12 关节角 + 2 夹爪）
    \item 删无效任务名：\texttt{unknown\_task} / \texttt{n/a} / \texttt{none} / 空白
    \item delta EE：按 FPS 换成每秒量；平移速度 $>0.5$ 或旋转速度 $>1.0$ 的步标记无效
    \item \textbf{按步 mask 而非丢整条轨迹}：chunk 内无效比例 $r_{\mathcal C} > 0.5$ 才丢弃；之后平移 / 旋转分别除以 0.5 / 1.0
    \item 关节角：先删 $[-2\pi, 2\pi]$ 之外的值，再 wrap 到 $[-\pi,\pi]$，不再额外归一化
    \item 夹爪：去极值后逐数据集 min-max 到 $[0,1]$
    \item 辅助数据：LLaVA-ReCap-CC3M、LLaVA-OneVision（caption）；RefCOCO、COCO-ReM（BBox-QA）；PixMo-Points、RoboPoint（Point-QA）；SenseNova-SI-800K（Spatial-QA）；Nemotron-SFT（纯文本）
  \end{itemize}
\end{frame}

\begin{frame}{Backup：VLAct 真机设置与结果（附录 J）}
  \begin{columns}[T]
    \column{0.46\textwidth}
    \textbf{设置}
    \begin{itemize}
      \item Franka Research 3 7-DoF（单臂 / 双臂各一模型）
      \item 外部 RealSense D435 + 腕部 D405，图像 224$\times$224
      \item GELLO 遥操作采集；单臂任务 50 条、双臂 100 条示教
      \item 50k 步，8$\times$H800；每任务 10 次 rollout，初始状态固定共享
      \item 长程任务按步计分（table cleaning 3 步各 0.33）
    \end{itemize}
    \column{0.50\textwidth}
    \begin{center}\footnotesize
    \begin{tabular}{p{3.6cm}rr}
      \toprule
      任务组 & VLAct & 基线 OFT \\
      \midrule
      单臂短程 ID（4 任务） & 92.5 & 77.5 \\
      短程 OOD：pot 新物体 & 90.0 & 73.3 \\
      短程 OOD：cup 新物体 & 90.0 & 65.0 \\
      长程：table cleaning & 86.6 & 73.3 \\
      长程：scoop beans & 80.0 & 33.3 \\
      长程 OOD：扩展 / 全替换 & 82.5 / 83.3 & 47.5 / 46.6 \\
      双臂协调（5 任务） & 72.0 & 44.0 \\
      \bottomrule
    \end{tabular}
    \end{center}
    \footnotesize 基线常见失败：立方体位置超出训练分布时冻住；scoop beans 跳过舀取直接空勺倒；叠毛巾双臂速度不匹配导致滑落。
  \end{columns}
\end{frame}

\begin{frame}{Backup：动作头选型指南}
  \begin{center}\small
  \begin{tabular}{p{5.2cm}p{2.6cm}p{5.4cm}}
    \toprule
    场景 & 建议 & 理由 \\
    \midrule
    单本体、数据充足、求简单 & OFT & 与生成式头持平，一次前向 \\
    低数据（$\le$50 条/任务）或强扰动 & OFT & RoboTwin Base/Random 领先 18 点 \\
    需要多模态动作分布 & PI / GR00T & 回归头是点估计 \\
    需要状态条件、独立运动模块 & GR00T & 状态与本体标识进 System 1 \\
    只能改 LLM 输出层 & FAST & 唯一不改架构，接受 15$\sim$30 点代价 \\
    做持续预训练供他人复用 & \textbf{$\ge$2 个连续头共监督} & 单头造成 decoder lock-in \\
    多本体训练 & 零填充 + 夹爪对齐 + mask & $\alpha$ 表 6 与 VLAct 表 12 叠加 \\
    关节角控制 & wrap + wrap loss & +5 点，零成本 \\
    \bottomrule
  \end{tabular}
  \end{center}
\end{frame}

\end{document}
